\documentclass[varwidth=\maxdimen,border=4pt]{standalone}
\usepackage{booktabs,array,fontspec}
\usepackage{xeCJK}
\setmainfont{Times New Roman}
\setCJKmainfont{Apple SD Gothic Neo}
\xeCJKsetup{CJKspace=true}
\renewcommand{\arraystretch}{1.12}
\begin{document}\footnotesize
\newsavebox{\tb}\sbox{\tb}{%
\begin{tabular}{l l p{5.2cm} p{4.2cm} l p{3.6cm}}
\toprule
Construct & Variable (by wave) & Questionnaire wording (verbatim) & English label & Block & Role \\
\midrule
IA3 & Q11\_3/문11\_3 & 중도포기 또는 중단된 혁신활동 & innovation activity abandoned or suspended & 문11 & Outcome: abandoned/suspended innovation activity (quit) \\
PROT8 & Q2\_8/문2\_8 & 경쟁사에 앞서 시장 선점 & lead-time advantage: entering the market ahead of competitors & 문2 & Outcome: lead-time protection (pre-emption) \\
ME4 & Q1K2\_4/문1\_2\_4 & 새로운 경쟁사의 진입이 귀사의 시장 지위를 위협함 & entry of new competitors threatens market position & 문1-2 & Treatment: perceived entry threat \\
RDSTAFF & Q33K2\_1\_3/문33\_4\_1 & 상용 근로자 중 연구개발(R\&D) 전담인력 비율(\%) & share of dedicated R\&D staff among regular employees (\%) & 문33 & Gap proxy: dedicated R\&D staff share \\
ME1 & Q1K2\_1/문1\_2\_1 & 상품(제품 또는 서비스)이 빠르게 구식이 되어버림 & products become outdated quickly & 문1-2 & Response-propensity index (me\_mean8) \\
ME2 & Q1K2\_2/문1\_2\_2 & 미래 기술 발전을 예측하기 어려움 & future technological developments hard to predict & 문1-2 & Index; control (tech. uncertainty) \\
ME3 & Q1K2\_3/문1\_2\_3 & 경쟁사 상품들이 귀사의 상품들을 쉽게 대체함 & competitor products easily substitute ours & 문1-2 & Index \\
ME5 & Q1K2\_5/문1\_2\_5 & 경쟁사의 행동을 예측하기 어려움 & competitors' actions hard to predict & 문1-2 & Index; exploratory sigma item \\
ME6 & Q1K2\_6/문1\_2\_6 & 수요의 변화를 예측하기 어려움 & changes in demand hard to predict & 문1-2 & Index; placebo \\
ME7 & Q1K2\_7/문1\_2\_7 & 해외 경쟁사들과의 경쟁이 치열함 & strong competition from abroad & 문1-2 & Index \\
ME8 & Q1K2\_8/문1\_2\_8 & 가격 인상이 고객의 손실을 초래함 & price increases lead to loss of customers & 문1-2 & Index; placebo \\
IA1 & Q11\_1/문11\_1 & 지난 3년간(2021\textasciitilde{}2023년) 기간 내 완료된 혁신활동 (상품혁신활동 또는 비즈니스프로세스혁신활동) & innovation activity completed within the period & 문11 & Sample filter; portfolio breadth \\
IA2 & Q11\_2/문11\_2 & 2023년 말 기준, 계속 진행 중이었던 혁신활동 & innovation activity ongoing at year-end & 문11 & Portfolio breadth \\
IA4 & Q11\_4/문11\_4 & 귀사 내부에서 전체를 독자적으로 자체수행한 연구개발(R\&D) 활동 & in-house R\&D & 문11 & Portfolio breadth \\
IA5 & Q11\_5/문11\_5 & 귀사와 타 기업 또는 타 기관이 계약을 통해 공동으로 협력하여 수행한 연구개발(R\&D) 활동 & collaborative R\&D & 문11 & Portfolio breadth \\
IA6 & Q11\_6/문11\_6 & 외주계약을 통해 타 기업 또는 타 기관에게 용역 전체를 의뢰하여 수행된 연구개발(R\&D) 활동 & contracted-out R\&D & 문11 & Portfolio breadth \\
EMP & Q33K1\_3/문33\_1\_3 & 2023년 상용 근로자 수 & regular employees in the reference year & 문33 & Control: regular employees (log) \\
FORM & Q31K1/문31\_1 & 기업 형태 & firm type: independent / domestic-group affiliate / foreign-group affiliate & 문31 & Control: group affiliation \\
LISTED & Q31K4/문31\_4 & 상장여부 & stock-market listing: KOSPI / KOSDAQ / KONEX / not listed & 문31 & Control: listing \\
MF1 & Q6\_1/문6\_1 & (시장최초 상품혁신) 경쟁사보다 앞서 기존에 없던 새롭거나 획기적으로 개선된 상품을 출시 & new-to-market product innovation (ahead of competitors) & 문6 & Exploratory outcome: new-to-market product \\
PI1 & Q5\_1/문5\_1 & (제품혁신) 귀사의 기존 제품 대비 새롭거나 획기적으로 개선된 제품 & product (goods) innovation introduced & 문5 & Filter for MF1 \\
PI2 & Q5\_2/문5\_2 & (서비스혁신) 귀사의 기존 서비스 대비 새롭거나 획기적으로 개선된 서비스 & service innovation introduced & 문5 & Filter for MF1 \\
\bottomrule
\end{tabular}}
\usebox{\tb}
\par\vspace{3pt}\parbox{\wd\tb}{\footnotesize Wording from the printed questionnaires of the four waves (identical across waves except reference years); variable names as in the STEPI provision files (2022/2023 문-style, 2024/2025 Q-style). Full per-wave item maps with printed codes and observed values: item\_maps/item\_map\_<wave>.json.}
\end{document}
